\documentclass[11pt]{article}
\usepackage[utf8]{inputenc}
\usepackage[T1]{fontenc}
\usepackage{amsmath,amssymb,mathtools}
\usepackage[margin=1in]{geometry}
\usepackage{hyperref}
\usepackage{array}
\usepackage{parskip}
\setlength{\parindent}{0pt}
\hypersetup{colorlinks=true,linkcolor=blue,urlcolor=blue}
\title{Lecture 21 — System Analysis Using the Unilateral Z-Transform}
\author{}
\date{}
\begin{document}
\maketitle

\textbf{Course:} CEC 315 — Signals and Systems (Spring 2026)
\textbf{Instructor:} Rogelio Gracia Otalvaro
\textbf{Source PDF:} \texttt{all\_lectures/cec315-lctr21-system-analysis-unilateral-z.pdf}
\textbf{Exam coverage:} Exam 3

\section*{Lctr 21: DT System Analysis and the Unilateral z-Transform}

Rogelio Gracia Otalvaro

Spring 2026

\textbf{Focus: Sections 10.7–10.9 (Pages 774–796)}

\section*{21.1 The Big Picture}

\subsection*{Why This Matters}

This lecture completes the z-transform trilogy by connecting the machinery to DT system analysis. The goals mirror Lecture 18 exactly:

\textbf{Part I:} Given a difference equation, find the system function $H(z)$, determine stability and causality, compute the output.

\textbf{Part II:} Use the unilateral z-transform to solve difference equations with nonzero initial conditions (the DT analog of solving IVPs with the unilateral Laplace transform).

After this lecture, you have a complete toolkit for both CT and DT system analysis.

\textbf{Roadmap:}

\begin{enumerate}
  \item The system function $H(z)$ from difference equations.
  \item Causality and stability in the z-domain.
  \item Complete system analysis example.
  \item Block diagram representations.
  \item The unilateral z-transform and its time-shift property.
  \item Solving difference equations with initial conditions.
\end{enumerate}

\section*{21.2 The System Function H(z)}

A causal LTI DT system described by a linear constant-coefficient difference equation:

\[\sum_{k=0}^{N} a_k\, y[n-k] = \sum_{k=0}^{M} b_k\, x[n-k]\]

Taking the bilateral z-transform (zero initial conditions): each delay $y[n-k]$ becomes $z^{-k} Y(z)$, each $x[n-k]$ becomes $z^{-k} X(z)$:

\[\left(\sum_{k=0}^{N} a_k z^{-k}\right) Y(z) = \left(\sum_{k=0}^{M} b_k z^{-k}\right) X(z)\]

\[H(z) = \frac{Y(z)}{X(z)} = \frac{\sum_{k=0}^{M} b_k z^{-k}}{\sum_{k=0}^{N} a_k z^{-k}} = \frac{b_0 + b_1 z^{-1} + \cdots + b_M z^{-M}}{a_0 + a_1 z^{-1} + \cdots + a_N z^{-N}}\]

\subsection*{Key Insight}

The z-transform converts the difference equation into an algebraic equation: each unit delay $n \to n-1$ becomes multiplication by $z^{-1}$. This is why $z^{-1}$ is called the "delay operator" in digital signal processing.

\subsection*{Example: From Difference Equation to H(z)}

\textbf{System:} $y[n] - 0.8\, y[n-1] + 0.15\, y[n-2] = x[n] + 2x[n-1]$

\textbf{Step 1:} Replace delays with $z^{-k}$:

\[(1 - 0.8 z^{-1} + 0.15 z^{-2}) Y(z) = (1 + 2 z^{-1}) X(z)\]

\textbf{Step 2:} Form $H(z)$:

\[H(z) = \frac{1 + 2 z^{-1}}{1 - 0.8 z^{-1} + 0.15 z^{-2}}\]

\textbf{Step 3:} Factor the denominator. We need roots of $1 - 0.8 z^{-1} + 0.15 z^{-2} = 0$, i.e., $z^2 - 0.8 z + 0.15 = 0$:

\[z = \frac{0.8 \pm \sqrt{0.64 - 0.60}}{2} = \frac{0.8 \pm 0.2}{2}\]

$z_1 = 0.5$, $z_2 = 0.3$.

So:

\[H(z) = \frac{1 + 2 z^{-1}}{(1 - 0.5 z^{-1})(1 - 0.3 z^{-1})}\]

\textbf{Poles:} $z = 0.5, 0.3$ (both inside the unit circle). \textbf{Zero:} from $1 + 2 z^{-1} = 0$, $z = -2$.

\section*{21.3 Causality and Stability}

\subsection*{Causality Criterion}

A DT LTI system is causal iff $h[n] = 0$ for $n < 0$, which requires the ROC to be the exterior of a circle: $|z| > r_{\max}$ (outside the outermost pole).

For rational $H(z)$: causality also requires $M \leq N$ (numerator degree $\leq$ denominator degree in $z^{-1}$), or equivalently $H(z) \to$ finite as $z \to \infty$.

\subsection*{Stability Criterion}

A DT LTI system is BIBO stable iff $\sum |h[n]| < \infty$, which requires the ROC to include the unit circle ($|z| = 1$).

\subsection*{The Golden Rule for Causal DT Systems}

A causal LTI DT system with rational $H(z)$ is \textbf{BIBO stable if and only if all poles lie strictly inside the unit circle}:

\[\text{Causal + Stable} \iff |p_i| < 1 \ \forall \ \text{poles}\ p_i\]

\textbf{Why:} For a causal system, ROC is $|z| > |p_{\max}|$. For the unit circle to be in the ROC, we need $|p_{\max}| < 1$.

\textit{Figure: Two pole-zero plots in the z-plane. Left — "Causal + Stable": all poles ($\times$) lie inside the unit circle $|z| = 1$. Right — "Causal + Unstable": at least one pole lies outside the unit circle.}

\subsection*{Stability Classification}

\begin{center}
\begin{tabular}{|l|l|l|l|}
\hline
\textbf{$H(z)$} & \textbf{Poles} & \textbf{$\|p_i\|$} & \textbf{Causal + Stable?} \\
\hline
$\dfrac{1}{1 - 0.5 z^{-1}}$ & $z = 0.5$ & $0.5 < 1$ & Yes \\
\hline
$\dfrac{1}{1 - 1.2 z^{-1}}$ & $z = 1.2$ & $1.2 > 1$ & No (unstable) \\
\hline
$\dfrac{1}{1 - z^{-1}}$ & $z = 1$ & $1 = 1$ & Marginally stable \\
\hline
$\dfrac{1}{(1 - 0.3 z^{-1})(1 + 0.7 z^{-1})}$ & $z = 0.3, -0.7$ & $0.3, 0.7 < 1$ & Yes \\
\hline
\end{tabular}
\end{center}

\section*{21.4 Complete System Analysis Pipeline}

\subsection*{Full Pipeline: Difference Equation to Output}

\textbf{System:} $y[n] - 0.5\, y[n-1] = x[n]$, \textbf{Input:} $x[n] = (0.8)^n u[n]$.

\textbf{Step 1:} Find $H(z)$.

\[(1 - 0.5 z^{-1}) Y = X \implies H(z) = \frac{1}{1 - 0.5 z^{-1}}, \quad |z| > 0.5\]

Pole at $z = 0.5$ (inside unit circle $\Rightarrow$ causal and stable). \checkmark

\textbf{Step 2:} Find $X(z)$.

\[X(z) = \frac{1}{1 - 0.8 z^{-1}}, \quad |z| > 0.8\]

\textbf{Step 3:} Find $Y(z) = H(z) \cdot X(z)$.

\[Y(z) = \frac{1}{(1 - 0.5 z^{-1})(1 - 0.8 z^{-1})}, \quad |z| > 0.8\]

\textbf{Step 4:} Partial fractions.

\[\frac{1}{(1 - 0.5 z^{-1})(1 - 0.8 z^{-1})} = \frac{A}{1 - 0.5 z^{-1}} + \frac{B}{1 - 0.8 z^{-1}}\]

Set $z^{-1} = 2$: $1 = A(1 - 1.6) \Rightarrow A = -5/3$.
Set $z^{-1} = 1.25$: $1 = B(1 - 0.625) \Rightarrow B = 8/3$.

\textbf{Step 5:} Invert. Both poles inside ROC $\Rightarrow$ both right-sided:

\[y[n] = \left[-\frac{5}{3}(0.5)^n + \frac{8}{3}(0.8)^n\right] u[n]\]

\textbf{Step 6:} Verify.

\begin{itemize}
  \item $y[0] = -5/3 + 8/3 = 3/3 = 1 = x[0]$ \checkmark (from the difference equation with $y[-1] = 0$)
  \item $y[1] = -5/3 \cdot 0.5 + 8/3 \cdot 0.8 = -5/6 + 6.4/3 = -5/6 + 32/15$. From the ODE: $y[1] = 0.5 y[0] + x[1] = 0.5 + 0.8 = 1.3$. Check: $-5/6 + 32/15 = -25/30 + 64/30 = 39/30 = 1.3$ \checkmark
\end{itemize}

\section*{21.5 Block Diagram Representations}

In the z-domain, the three basic building blocks are:

\textit{Figure: Three basic z-domain building blocks — (1) Unit Delay: input $X(z)$ through a $z^{-1}$ block outputs $z^{-1} X(z)$; (2) Gain: a multiplier block labeled $a$; (3) Adder: a summing junction marked with $+$.}

Every rational $H(z)$ can be implemented as a combination of delays ($z^{-1}$), multipliers (gains), and adders. This is the basis of digital filter design.

\subsection*{Key Insight}

The block $z^{-1}$ represents storing one sample in memory. A difference equation of order $N$ requires exactly $N$ memory elements. In software, each $z^{-1}$ is one element of a circular buffer; in hardware, it is one register or latch. This direct correspondence between the z-domain and digital implementation is why the z-transform is the language of DSP.

\section*{21.6 Definition}

\subsection*{The Unilateral z-Transform}

\[\mathcal{X}(z) = \sum_{n=0}^{\infty} x[n]\, z^{-n}\]

The sum starts at $n = 0$, not $n = -\infty$. For causal signals ($x[n] = 0$ for $n < 0$), the unilateral and bilateral transforms are identical.

\section*{21.7 The Key Property: Time Shifting}

The bilateral z-transform gives: $x[n-1] \xleftrightarrow{\mathcal{Z}} z^{-1} X(z)$.

The unilateral z-transform gives:

\[x[n-1] \xleftrightarrow{\mathcal{Z}_u} z^{-1} \mathcal{X}(z) + x[-1]\]

\[x[n-2] \xleftrightarrow{\mathcal{Z}_u} z^{-2} \mathcal{X}(z) + x[-2] + x[-1]\, z^{-1}\]

General pattern for delay by $m$:

\[x[n-m] \xleftrightarrow{\mathcal{Z}_u} z^{-m} \mathcal{X}(z) + x[-m] + x[-m+1]\, z^{-1} + \cdots + x[-1]\, z^{-(m-1)}\]

\subsection*{Key Insight}

The initial conditions $x[-1], x[-2], \ldots$ appear explicitly, just as $x(0^-), x'(0^-)$ appeared in the unilateral Laplace differentiation property. This is the mechanism for incorporating nonzero initial conditions into difference equation solutions.

\section*{21.8 Solving Difference Equations with Initial Conditions}

\subsection*{First-Order Difference Equation with IC}

\textbf{Solve:} $y[n] - 0.6\, y[n-1] = (0.5)^n u[n]$, $\ y[-1] = 4$.

\textbf{Step 1:} Unilateral z-transform.

For the left side:

\[Y(z) - 0.6\left(z^{-1} Y(z) + y[-1]\right) = Y(z) - 0.6 z^{-1} Y(z) - 0.6 \cdot 4\]

\[= (1 - 0.6 z^{-1}) Y(z) - 2.4\]

For the right side: $\dfrac{1}{1 - 0.5 z^{-1}}$.

\textbf{Step 2:} Solve for $Y(z)$.

\[(1 - 0.6 z^{-1}) Y(z) = \frac{1}{1 - 0.5 z^{-1}} + 2.4\]

\[Y(z) = \frac{1}{(1 - 0.5 z^{-1})(1 - 0.6 z^{-1})} + \frac{2.4}{1 - 0.6 z^{-1}}\]

\textbf{Step 3:} Partial fractions on the first term.

\[\frac{1}{(1 - 0.5 z^{-1})(1 - 0.6 z^{-1})} = \frac{A}{1 - 0.5 z^{-1}} + \frac{B}{1 - 0.6 z^{-1}}\]

Set $z^{-1} = 2$: $1 = A(1 - 1.2) \Rightarrow A = -5$.
Set $z^{-1} = 5/3$: $1 = B(1 - 5/6) = B/6 \Rightarrow B = 6$.

So:

\[Y(z) = \frac{-5}{1 - 0.5 z^{-1}} + \frac{6}{1 - 0.6 z^{-1}} + \frac{2.4}{1 - 0.6 z^{-1}} = \frac{-5}{1 - 0.5 z^{-1}} + \frac{8.4}{1 - 0.6 z^{-1}}\]

\textbf{Step 4:} Invert.

\[y[n] = \left[-5 \cdot (0.5)^n + 8.4 \cdot (0.6)^n\right] u[n]\]

\textbf{Step 5:} Verify.

\begin{itemize}
  \item $y[0] = -5 + 8.4 = 3.4$. From the ODE: $y[0] = 0.6 \cdot y[-1] + x[0] = 0.6 \cdot 4 + 1 = 3.4$ \checkmark.
  \item $y[1] = -5(0.5) + 8.4(0.6) = -2.5 + 5.04 = 2.54$. From the ODE: $y[1] = 0.6 \cdot y[0] + x[1] = 0.6 \cdot 3.4 + 0.5 = 2.54$ \checkmark.
\end{itemize}

\subsection*{Second-Order}

\textbf{Solve:} $y[n] - 0.7\, y[n-1] + 0.1\, y[n-2] = 0$, $\ y[-1] = 1, \ y[-2] = 0$.

\textbf{Step 1:} Unilateral z-transform.

\[\mathcal{Z}_u\{y[n-1]\} = z^{-1} Y(z) + y[-1] = z^{-1} Y + 1\]

\[\mathcal{Z}_u\{y[n-2]\} = z^{-2} Y(z) + y[-2] + y[-1] z^{-1} = z^{-2} Y + 0 + z^{-1}\]

Substitute:

\[Y - 0.7(z^{-1} Y + 1) + 0.1(z^{-2} Y + z^{-1}) = 0\]

\[(1 - 0.7 z^{-1} + 0.1 z^{-2}) Y = 0.7 - 0.1 z^{-1}\]

\textbf{Step 2:} Factor. $z^2 - 0.7 z + 0.1 = 0 \Rightarrow z = \dfrac{0.7 \pm \sqrt{0.49 - 0.4}}{2} = \dfrac{0.7 \pm 0.3}{2}$. Poles: $z = 0.5$ and $z = 0.2$.

\[Y(z) = \frac{0.7 - 0.1 z^{-1}}{(1 - 0.5 z^{-1})(1 - 0.2 z^{-1})}\]

\textbf{Step 3:} Partial fractions.

\[\frac{0.7 - 0.1 z^{-1}}{(1 - 0.5 z^{-1})(1 - 0.2 z^{-1})} = \frac{A}{1 - 0.5 z^{-1}} + \frac{B}{1 - 0.2 z^{-1}}\]

Set $z^{-1} = 2$: $0.7 - 0.2 = A(1 - 0.4) \Rightarrow 0.5 = 0.6 A \Rightarrow A = 5/6$.
Set $z^{-1} = 5$: $0.7 - 0.5 = B(1 - 2.5) \Rightarrow 0.2 = -1.5 B \Rightarrow B = -2/15$.

\textbf{Step 4:} Invert.

\[y[n] = \left[\frac{5}{6} \cdot (0.5)^n - \frac{2}{15} \cdot (0.2)^n\right] u[n]\]

This is a \textbf{pure zero-input response}: there is no external input, so the output is entirely due to the initial conditions. Both modes decay because both poles are inside the unit circle.

\textbf{Verify:} $y[0] = 5/6 - 2/15 = 25/30 - 4/30 = 21/30 = 0.7$. From ODE: $y[0] = 0.7 \cdot y[-1] - 0.1 \cdot y[-2] = 0.7 \cdot 1 - 0 = 0.7$ \checkmark.

\section*{21.9 ZSR + ZIR Decomposition}

Just as in CT (Lecture 18):

\[y[n] = \underbrace{y_{ZS}[n]}_{\text{zero-state}} + \underbrace{y_{ZI}[n]}_{\text{zero-input}}\]

$y_{ZS}[n]$: output due to input alone, with all initial conditions set to zero.
$y_{ZI}[n]$: output due to initial conditions alone, with zero input.

\section*{21.10 Summary of the Three-Lecture z-Transform Block}

\textit{Figure: Flow diagram showing the z-transform progression — Lecture 19 (Definition, ROC) $\to$ [tools] $\to$ Lecture 20 (Inverse, Properties) $\to$ [application] $\to$ Lecture 21 (Systems, IVPs).}

\begin{center}
\begin{tabular}{|l|l|l|}
\hline
\textbf{} & \textbf{Laplace (Lctrs 16–18)} & \textbf{z-Transform (Lctrs 19–21)} \\
\hline
Converts & differential eqs $\to$ algebra & difference eqs $\to$ algebra \\
\hline
Delay operator & $e^{-s t_0}$ & $z^{-1}$ \\
\hline
Causal + Stable & poles in LHP ($\mathrm{Re}\{p\} < 0$) & poles inside $\|z\| = 1$ ($\|p\| < 1$) \\
\hline
Unilateral IC property & $y'(t) \to sY - y(0^-)$ & $y[n-1] \to z^{-1} Y + y[-1]$ \\
\hline
FT exists & $j\omega$-axis in ROC & unit circle in ROC \\
\hline
\end{tabular}
\end{center}

\textbf{Key takeaways:}

\begin{enumerate}
  \item $H(z)$ from a difference equation: replace each delay $n - k$ with $z^{-k}$ and solve for $Y/X$.
  \item Causal + stable $\iff$ all poles strictly inside the unit circle.
  \item Block diagrams use $z^{-1}$ as the basic delay element; this maps directly to digital hardware/software.
  \item The unilateral z-transform handles initial conditions via the modified time-shift property: $y[n-1] \to z^{-1} Y + y[-1]$.
  \item Total response = ZSR + ZIR.
\end{enumerate}

\section*{21.11 Common Mistakes to Avoid}

\begin{enumerate}
  \item \textbf{Confusing "inside unit circle" with "LHP":} In the z-domain, stability requires poles inside the unit circle (not in the left half-plane). A pole at $z = -0.9$ is stable (inside the circle) even though it is on the negative real axis.
  \item \textbf{Getting the unilateral shift property wrong:} For $y[n-1]$: the IC term is $+y[-1]$, not $-y[-1]$. Compare with Laplace: $y'(t) \to sY - y(0^-)$ has a minus sign, but the z-domain version has a plus sign.
  \item \textbf{Forgetting that $z^{-1}$ means delay, not advance:} $x[n-1] \leftrightarrow z^{-1} X(z)$ is a delay. $x[n+1] \leftrightarrow z X(z) - z x[0]$ is an advance (unilateral). Do not confuse the two.
  \item \textbf{Misidentifying pole locations from the $z^{-1}$ form:} In $H(z) = \dfrac{1}{1 - a z^{-1}}$, the pole is at $z = a$, not $z = 1/a$ or $z = -a$.
  \item \textbf{Checking stability by testing $a < 1$ instead of $|a| < 1$:} A pole at $z = -0.9$ has $|z| = 0.9 < 1$, so it is stable. But $a = -0.9 < 1$ is not the right check; $|a| < 1$ is.
\end{enumerate}

\textit{Rogelio Gracia Otalvaro}

\section*{Practice Problems Summary}

\begin{itemize}
  \item \textbf{Example (Section 21.2) — From Difference Equation to H(z):} Given $y[n] - 0.8 y[n-1] + 0.15 y[n-2] = x[n] + 2 x[n-1]$, derive $H(z)$, factor, and identify poles $z = 0.5, 0.3$ and zero $z = -2$.
  \item \textbf{Stability Classification Table (Section 21.3):} Classify four sample transfer functions as causal-and-stable, unstable, or marginally stable based on pole locations relative to the unit circle.
  \item \textbf{Full Pipeline Example (Section 21.4):} Solve $y[n] - 0.5 y[n-1] = x[n]$ with $x[n] = (0.8)^n u[n]$ end-to-end: find $H(z)$, $X(z)$, $Y(z)$, partial fractions, inversion, and numerical verification.
  \item \textbf{First-Order Difference Equation with IC (Section 21.8):} Solve $y[n] - 0.6 y[n-1] = (0.5)^n u[n]$ with $y[-1] = 4$ using the unilateral z-transform, partial fractions, and verify $y[0], y[1]$.
  \item \textbf{Second-Order Zero-Input Response (Section 21.8):} Solve the homogeneous $y[n] - 0.7 y[n-1] + 0.1 y[n-2] = 0$ with $y[-1] = 1, y[-2] = 0$, demonstrating a pure zero-input response.
\end{itemize}

\end{document}
